\worldchapter{Wunden und Erholung}{wounds}
\label{ch:wunden}

\begin{multicols*}{2}
\section{Wunden}\label{sec:wunden}\index{Wunde}\index{Wundengrenze}

\textbf{Wunden} stellen körperliche Verletzungen dar.
Jeglicher Schaden wird als Wunde festgehalten.

Die Wundschwelle eines Charakters gibt an, ab welcher Anzahl an Wunden er als \textbf{Sterbend} gilt.

\[
\text{Wundschwelle} = 3 + \text{Körper}
\]

Es gilt stets nur die höchste aktuelle Wundstufe:

\begin{description}[style=nextline,leftmargin=0pt]
\item[\textbf{Verletzt} (1+)] die erste körperliche Probe jeder Runde ist -10
\item[\textbf{Beeinträchtigt} (3+)] alle körperlichen Proben -10, kein Sprint
\item[\textbf{Kritisch} (5+)] alle Proben -20, beim Zug 1 Aktion weniger auffrischen, Minimum 1
\item[\textbf{Sterbend} (> Threshold)] am Ende jedes eigenen Zuges +1 Wunde
\end{description}

Wenn die Wunden die Wundschwelle um 3 übersteigen, stirbt der Charakter.

\section{Stabilisieren}\label{sec:stabilisieren}\index{Stabilisieren}

Ein naher oder gebundener Verbündeter kann \textbf{Stabilisieren} als Aktion versuchen:

\[
30 + \text{Geist oder Wille} + \text{Heilen}
\]

Bei Erfolg wird die Anzahl der Wunden exakt auf die Wundschwelle gesetzt.
\textbf{Sterbend} und \textbf{Blutend} werden beendet.

Bei knappem oder klarem Misserfolg: keine Änderung.

Bei schwerem oder katastrophalem Misserfolg: keine Änderung und sofort +1 Wunde.

\section{Bürde}\label{sec:burde}\index{Bürde}

\textbf{Bürde} sammelt Stress, Erschöpfung, Schuld, moralischen Druck und Verderbnislast.
Die Bürdenschwelle verdeutlicht, wie viel Last ein Charakter aushalten kann.

\[
\text{Bürdenschwelle} = 5 + \lfloor \text{Wille} / 2 \rfloor
\]

\begin{description}[style=nextline,leftmargin=0pt]
\item[\textbf{Stabil}] 0--2
\item[\textbf{Belastet}] 3 bis Bürdenschwelle - 1
\item[\textbf{Wendepunkt}]\index{Wendepunkt} exakt Bürdenschwelle
\end{description}

Am Wendepunkt wird entweder sofort ein neues Mal genommen oder der Charakter bricht zusammen und verliert für den aktuellen Kampf eine Aktion pro Kampfrunde.

\section{Verderbnis}\label{sec:verderbnis}\index{Verderbnis}

\textbf{Verderbnis} ist die Langzeitfolge verbotener Praktiken.
%TODO: Define this!

\begin{itemize}[leftmargin=*]
\item \textbf{0--1}: noch verborgen
\item \textbf{2}: chaotische Zauberei gilt immer
\item \textbf{3}: sofort neues Mal oder permanentes Tabu wählen, dann Verderbnis auf 1 setzen
\end{itemize}

\section{Permanentes Tabu}\label{sec:permanentes-tabu}

Ein \textbf{Tabu}\index{Tabu} ist ein niedergeschriebener verbotener Akt.
Tabus können von Abstammung oder Weg kommen, oder aber auch selbst auferlegt sein.

\begin{itemize}[leftmargin=*]
\item Wird das Tabu gebrochen: sofort -1 Omen
\item falls kein Omen vorhanden: stattdessen +1 Bürde
\item bis zu einer Sühneszene wird keine Gunst regeneriert
\end{itemize}

\section{Zustände}\label{sec:zustande}\index{Zustand}

Ein Charakter kann mehrere Zustände haben.
Diese ergeben sich aus der Situation heraus, so kann Feuer den Zustand \textbf{Brennend} verursachen.
Waffen können eine Eigenschaft haben, die bei einem Treffer einen Zustand verursacht.

\begin{description}[style=nextline,leftmargin=0pt]
\item[\textbf{Blutend}]\index{Blutend} 1 Wunde am Ende jedes eigenen Zuges, bis behandelt
\item[\textbf{Brennend}]\index{Brennend} 1 Wunde, wenn der Charakter handelt; endet durch Ersticken oder Löschen
\item[\textbf{Erschüttert}]\index{Erschüttert} -10 auf die nächste körperliche Probe; verschwindet danach oder durch Sich fassen/Unterstützen
\item[\textbf{Festgesetzt}]\index{Festgesetzt} keine Bewegung und kein Sprinten
\item[\textbf{Defekt}]\index{Defekt} die ausgewählte Ausrüstung ist unbrauchbar, bis sie repariert wurde
\end{description}

\section{Feldreparatur}\index{Feldreparatur}

Eine kaputte Waffe, ein Schild, ein Rüstungsteil oder Werkzeug kann im Feld geflickt werden:

\[
30 + \text{Geist oder Geschick} + \text{Handwerk}
\]

Kosten: 1 Aktion, Ziel auf sich selbst oder \textbf{Nah}.

\begin{itemize}[leftmargin=*]
\item Kritischer, starker oder sauberer Erfolg: Gegenstand funktioniert für den Rest der Szene
\item Knapp geschafft: funktioniert, ist aber nur geflickt; bricht er erneut in derselben Szene, ist keine weitere Feldreparatur möglich.
\item Knapp oder klar gescheitert: keine Änderung
\item Schwer oder katastrophal gescheitert: keine Änderung und das benötigte Ersatzmaterial ist verbraucht
\end{itemize}

\section{Erholung}\label{sec:erholung}\index{Erholung}

Es gibt folgende Möglichkeiten der Erholung.

\begin{description}[style=nextline,leftmargin=0pt]
\item[\textbf{Feldverband} (1--2 Minuten)]\index{Feldverband} stoppt Bluten, heilt aber keine Wunden
\item[\textbf{Kurze Rast} (10--20 Minuten sicher)]\index{Kurze Rast} eine Heilung-Probe; bei Erfolg 1 Wunde oder 1 Bürde entfernen.
\item[\textbf{Durchatmen}]\index{Durchatmen} einmal pro Szene 1 Bürde oder 1 Arkana oder 1 Gunst zurückgewinnen
\item[\textbf{Vollständige Rast}]\index{Vollständige Rast} 2 Punkte insgesamt aus Wunden und Bürde löschen, Arkana und Gunst voll auffüllen
\item[\textbf{Unsichere Rast}]\index{Unsichere Rast} nur 1 Punkt aus Wunden oder Bürde löschen; Arkana und Gunst bis zur Hälfte auffüllen
\end{description}

Ruhe entfernt niemals ein Mal.
\end{multicols*}
